% ============================================================
%  MTH636 – Game Theory | Lecture 19
%  Normal Form Representation \& Subgame Perfect Equilibrium
% ============================================================
\documentclass[12pt, a4paper]{article}

% ============================================================
%  MTH636 -- Game Theory
%  Notes Preamble — Color-Coded, TikZ-Enhanced
% ============================================================

% --- Core packages ---
\usepackage[T1]{fontenc}
\usepackage[utf8]{inputenc}
\usepackage{lmodern}
\usepackage[margin=2.5cm, top=2.8cm, bottom=2.8cm]{geometry}
\usepackage{microtype}

% --- Math ---
\usepackage{amsmath, amssymb, amsthm, mathtools}
\usepackage{bm}        % bold math
\usepackage{cancel}    % strike-through in math

% --- Color & Boxes ---
\usepackage[dvipsnames, svgnames, x11names, table]{xcolor}
\usepackage[most]{tcolorbox}
\tcbuselibrary{breakable, skins, theorems}

% --- TikZ & Diagrams ---
\usepackage{tikz}
\usepackage{tikz-qtree}
\usetikzlibrary{arrows.meta, positioning, shapes.geometric, fit, backgrounds,
                decorations.pathreplacing, calc, matrix, trees}
\usepackage{pgfplots}
\pgfplotsset{compat=1.18}

% --- Tables ---
\usepackage{booktabs}
\usepackage{multirow}
\usepackage{array}
\usepackage{tabularx}
\usepackage{makecell}

% --- Typography & Layout ---
\usepackage{enumitem}
\usepackage{parskip}
\usepackage{fancyhdr}
\usepackage{titlesec}
\usepackage{titletoc}
\usepackage{hyperref}
\usepackage{graphicx}
\usepackage{float}
\usepackage{caption}

% --- Color Palette ---
\definecolor{defblue}{HTML}{1A5276}      % Definition header
\definecolor{defbluebg}{HTML}{D6EAF8}    % Definition background
\definecolor{thmgreen}{HTML}{1E8449}     % Theorem header
\definecolor{thmbg}{HTML}{D5F5E3}        % Theorem background
\definecolor{exorange}{HTML}{B7770D}     % Example header
\definecolor{exbg}{HTML}{FEF9E7}         % Example background
\definecolor{warnred}{HTML}{922B21}      % Warning/Key Insight header
\definecolor{warnbg}{HTML}{FADBD8}       % Warning background
\definecolor{noticegray}{HTML}{2C3E50}   % Note header
\definecolor{noticebg}{HTML}{EBF5FB}     % Note background
\definecolor{taketawayyellow}{HTML}{7D6608} % Key Takeaway header
\definecolor{takeawaybg}{HTML}{FDFEFE}   % Takeaway background
\definecolor{accentpurple}{HTML}{6C3483} % Accent / proof
\definecolor{lightgray}{HTML}{F2F3F4}    % shading

% ============================================================
%  TCOLORBOX ENVIRONMENTS
% ============================================================

% -- Definition (Blue) --
\newtcolorbox{defbox}[2][]{
  enhanced, breakable,
  colback=defbluebg, colframe=defblue, coltitle=white,
  fonttitle=\bfseries\sffamily,
  title={Definition: #2},
  attach boxed title to top left={yshift=-2mm, xshift=4mm},
  boxed title style={colback=defblue, rounded corners},
  arc=4pt, boxrule=1.2pt,
  left=6pt, right=6pt, top=8pt, bottom=6pt,
  #1
}

% -- Theorem / Proposition (Green) --
\newtcolorbox{thmbox}[2][]{
  enhanced, breakable,
  colback=thmbg, colframe=thmgreen, coltitle=white,
  fonttitle=\bfseries\sffamily,
  title={Theorem: #2},
  attach boxed title to top left={yshift=-2mm, xshift=4mm},
  boxed title style={colback=thmgreen, rounded corners},
  arc=4pt, boxrule=1.2pt,
  left=6pt, right=6pt, top=8pt, bottom=6pt,
  #1
}

% -- Proposition (Green variant) --
\newtcolorbox{propbox}[2][]{
  enhanced, breakable,
  colback=thmbg, colframe=thmgreen!70!black, coltitle=white,
  fonttitle=\bfseries\sffamily,
  title={Proposition: #2},
  attach boxed title to top left={yshift=-2mm, xshift=4mm},
  boxed title style={colback=thmgreen!70!black, rounded corners},
  arc=4pt, boxrule=1.2pt,
  left=6pt, right=6pt, top=8pt, bottom=6pt,
  #1
}

% -- Example (Orange/Yellow) --
\newtcolorbox{exbox}[2][]{
  enhanced, breakable,
  colback=exbg, colframe=exorange, coltitle=white,
  fonttitle=\bfseries\sffamily,
  title={Example: #2},
  attach boxed title to top left={yshift=-2mm, xshift=4mm},
  boxed title style={colback=exorange, rounded corners},
  arc=4pt, boxrule=1.2pt,
  left=6pt, right=6pt, top=8pt, bottom=6pt,
  #1
}

% -- Key Insight / Warning (Red) --
\newtcolorbox{warnbox}[2][]{
  enhanced, breakable,
  colback=warnbg, colframe=warnred, coltitle=white,
  fonttitle=\bfseries\sffamily,
  title={Key Insight: #2},
  attach boxed title to top left={yshift=-2mm, xshift=4mm},
  boxed title style={colback=warnred, rounded corners},
  arc=4pt, boxrule=1.2pt,
  left=6pt, right=6pt, top=8pt, bottom=6pt,
  #1
}

% -- Note (Gray-Blue) --
\newtcolorbox{notebox}[1][]{
  enhanced, breakable,
  colback=noticebg, colframe=noticegray!60, coltitle=white,
  fonttitle=\bfseries\sffamily,
  title={Note},
  attach boxed title to top left={yshift=-2mm, xshift=4mm},
  boxed title style={colback=noticegray!60, rounded corners},
  arc=4pt, boxrule=0.8pt,
  left=6pt, right=6pt, top=8pt, bottom=6pt,
  #1
}

% -- Key Takeaway --
\newtcolorbox{takebox}[1][]{
  enhanced, breakable,
  colback=takeawaybg, colframe=taketawayyellow, coltitle=white,
  fonttitle=\bfseries\sffamily,
  title={Key Takeaway},
  attach boxed title to top left={yshift=-2mm, xshift=4mm},
  boxed title style={colback=taketawayyellow, rounded corners},
  arc=4pt, boxrule=1.2pt,
  left=6pt, right=6pt, top=8pt, bottom=6pt,
  #1
}

% -- Intuition --
\newtcolorbox{intubox}[1][]{
  enhanced, breakable,
  colback=accentpurple!8, colframe=accentpurple!60, coltitle=white,
  fonttitle=\bfseries\sffamily,
  title={Intuition},
  attach boxed title to top left={yshift=-2mm, xshift=4mm},
  boxed title style={colback=accentpurple!80, rounded corners},
  arc=4pt, boxrule=0.8pt,
  left=6pt, right=6pt, top=8pt, bottom=6pt,
  #1
}

% ============================================================
%  PAYOFF MATRIX HELPER
%  Usage: inside tabular, use \payoff{a}{b} for (a,b)
% ============================================================
\newcommand{\payoff}[2]{$(#1,\;#2)$}
\newcommand{\nash}[2]{\underline{$#1$},\underline{$#2$}}  % underline NE payoffs
\newcommand{\Nset}{\mathbb{N}}
\newcommand{\Rset}{\mathbb{R}}
\newcommand{\Sset}{\mathcal{S}}
\newcommand{\Ui}{u_i}

% ============================================================
%  SECTION FORMATTING
% ============================================================
\titleformat{\section}
  {\large\bfseries\sffamily\color{defblue}}
  {\thesection}{1em}{}[\titlerule]

\titleformat{\subsection}
  {\normalsize\bfseries\sffamily\color{thmgreen!80!black}}
  {\thesubsection}{1em}{}

\titleformat{\subsubsection}
  {\normalsize\itshape\sffamily\color{exorange!80!black}}
  {\thesubsubsection}{1em}{}

% ============================================================
%  HEADER / FOOTER
% ============================================================
\pagestyle{fancy}
\fancyhf{}
\fancyhead[L]{\sffamily\small\color{defblue} MTH636 -- Game Theory}
\fancyhead[R]{\sffamily\small\color{defblue} \leftmark}
\fancyfoot[C]{\sffamily\small\thepage}
\renewcommand{\headrulewidth}{0.4pt}

% ============================================================
%  HYPERREF SETUP
% ============================================================
\hypersetup{
  colorlinks=true,
  linkcolor=defblue,
  urlcolor=thmgreen,
  citecolor=accentpurple,
  bookmarksnumbered=true,
  pdftitle={MTH636 Game Theory Notes},
  pdfauthor={MTH636 Student Notes}
}

% ============================================================
%  MISC HELPERS
% ============================================================
\newcommand{\BR}{\text{BR}}        % Best Response
\newcommand{\NE}{\text{NE}}        % Nash Equilibrium
\newcommand{\SPNE}{\text{SPNE}}    % Subgame Perfect NE
\newcommand{\Dom}{\text{dom}}      % Dominance
\newcommand{\mixs}[1]{\sigma_{#1}} % Mixed strategy
\newcommand{\pures}[1]{s_{#1}}     % Pure strategy


% ============================================================
\begin{document}

% Title
\begin{center}
  {\LARGE\bfseries\sffamily\color{defblue} MTH636 -- Game Theory}\\[6pt]
  {\Large\sffamily Lecture 19: Normal Form equivalence \& SPE}\\[4pt]
\end{center}
\vspace{0.5cm}
\hrule height 1.5pt
\vspace{0.3cm}

% ============================================================
\section{Normal Form Representation of Extensive Form Games}

Suppose $\Gamma$ is an extensive form game (potentially with chance nodes). We can always translate it into an equivalent normal form game:
\[ \Gamma = \langle N, V, E, n_0, (V_i)_{i \in N \cup \{0\}}, (p_n)_{n \in V_0}, (u_i)_{i \in N} \rangle \]

The equivalent normal form game is given by $\langle N, (S_i)_{i \in N}, (\hat{u}_i)_{i \in N} \rangle$, where:
\begin{itemize}
    \item $N$ is the same set of players.
    \item $S_i$ is the set of pure strategies for player $i$ derived from the extensive form game (recall $s_i : V_i \to \cup A(n)$).
    \item $\hat{u}_i(s)$ is the expected utility of the strategy profile $s = (s_1, \dots, s_n)$.
\end{itemize}

\subsection{Utility Computation}
If the game has \textbf{no chance nodes}, a strategy profile $s$ uniquely traces a path to exactly one leaf node $\pi$. In this case:
\[ \hat{u}_i(s) = u_i(\pi) \]

If the game \textbf{has chance nodes}, a strategy profile $s$ no longer leads to a single leaf. Instead, because Nature introduces randomness at $V_0$ nodes, the profile $s$ induces a probability distribution over several possible leaf nodes. 
In this case, utilities are calculated as the \textbf{expected} utility according to the probability distribution over the leaf nodes $\pi$.
For example, if player 1 plays $L$ and player 2 plays $T$, and a chance node assigns 50/50 probability to reaching leaf $n_3$ or $n_4$, then:
\[ \hat{u}_1(L, T) = \frac{1}{2}u_1(n_3) + \frac{1}{2}u_1(n_4) \]

\textbf{Note on Zermelo's Theorem:}
When chance nodes are introduced, Zermelo/Von Neumann's theorem (which guaranteed a pure winning or drawing strategy) generally \textit{does not hold}. This is because the introduction of expected utilities smooths the discrete $\{1, -1, 0\}$ outcomes into fractional expected payoffs, breaking the deterministic Win/Loss/Draw structure.

% ============================================================
\section{Nash Equilibrium vs Subgame Perfect Equilibrium}

Once translated into normal form, we can find the Nash Equilibria using the standard definition:
\begin{defbox}{Nash Equilibrium}
  A strategy profile $s^*$ is a Nash Equilibrium if:
  \[ u_i(s^*) \geq u_i(s_i, s_{-i}^*) \quad \forall i \in N, \ \forall s_i \in S_i \]
\end{defbox}

\subsection{The Problem with Normal Form NE}
Consider the Sequential Battle of the Sexes from Lecture 15:
\begin{itemize}
    \item Man chooses $F$ or $D$ first.
    \item Woman observes and chooses $F$ or $D$.
\end{itemize}
The normal form matrix has a Nash Equilibrium at $(D, DD)$ which yields payoffs $(1,2)$. 

However, let's analyze the strategy $DD$ ("I will play Dance if you play Football, and I will play Dance if you play Dance"). 
If the Man actually chose Football ($F$), the game moves to the subgame where the Woman must choose. Her payoffs are $(2,1)$ for choosing $F$ and $(0,0)$ for choosing $D$. 
Clearly, in that subgame, she would rationally choose $F$ to get 1 instead of 0. 

Thus, her strategy to play $D$ in response to $F$ is a \textbf{non-credible threat}. She is threatening to tank both of their payoffs to $(0,0)$ to scare the Man into playing $D$ initially. While $(D, DD)$ is mathematically a Nash Equilibrium of the whole game, it relies on irrational behavior in a specific subgame.

\subsection{Subgame Perfect Equilibrium (SPE)}
To eliminate non-credible threats, we refine our equilibrium concept for extensive form games.

\begin{defbox}{Subgame Perfect Equilibrium (SPE)}
  A strategy profile is a Subgame Perfect Equilibrium (SPE) if it constitutes a Nash Equilibrium in \textbf{every} subgame of the original game.
\end{defbox}

By demanding that players act rationally not just overall, but from \textit{every single decision node} (even ones off the equilibrium path), SPE effectively rules out empty threats. In the example above, $(D, DD)$ is an NE, but \textbf{not} an SPE. The profile $(F, FD)$ is both an NE and an SPE.

\end{document}
